\section{The chain of theta lifts}\label{sec:theta-chain}

Fix a large mesh point $x$ in Hypothesis~\ref{hyp:one-shift-arithmetic}.
Take the field $E$ of Proposition~\ref{prop:field-window}, sets $T_2,T$
as in Lemma~\ref{lem:weight-count}, and one shift $u_0\in H$.
Each of the $m$ lifts uses the top places $(0,0,t)$,
$t\in T$, and the thick places $(0,1,s)$, $s\in T_2+u_0$.

\subsection{The nodes and the highest weights}

Choose even $k,m$ and odd $t_r>0$ as in
\citep[Proposition~5.15]{adriazola-grh}, so that $k,m$ satisfy
\eqref{eq:arithmetic-ranks}. Put $b'=xb/Z$, with $b,Z$ as in
\S\ref{sec:weights}, and set $h_0=2\lceil x^{402}\rceil$.

\begin{remark}[Open spacing compatibility]\label{rem:spacing-compatibility}
The standard record in the released source uses
$T_O=\lceil x^{377}\rceil$ and requires
\[
                         4T_O^2+2t_r\le h_0.
\]
The present choice cannot satisfy this requirement: for $x\ge2$ the
left side is at least $4x^{754}$, whereas
$h_0\le2(x^{402}+1)<4x^{754}$. Thus the chain below is not presently
identified with that standard source record. Larger polynomial node
gaps remain compatible with the asymptotic ranges used in
\S\ref{sec:weights}, but this manuscript has not verified one common
choice together with the corresponding lifts, period maps and height
estimates. We therefore flag the incompatibility rather than claim a
fix.
\end{remark}

For each $t\in T$ choose integers $a_t,b_t$ in the intervals
\[
  xb\le a_t\le xb+\frac{xb}{kx^{13}},\qquad
  b'\le b_t\le(1+x^{-2})b'.
\]
Both intervals have length greater than one for large $x$.
For the initial parameters in \citep[Proposition~18.7]{adriazola-grh}
start with
\[
 b_1(\rho)=
 \begin{cases}
 b_t+\lfloor b'/x\rfloor, & \rho=(0,0,t),\ t\in T,\\
 \lceil b\rceil,           & \rho=(0,1,s),\ s\in T_2+u_0,\\
 \lceil b'/(xn)\rceil,    & \text{otherwise}.
 \end{cases}
\]
If $\sum_\rho b_1(\rho)$ is odd, increase $b_1$ by one at a single
thick place. Set
\begin{equation}\label{eq:chain-nodes}
    w_j(\rho)=b_1(\rho)+(j-\tfrac12)h_0,\qquad 1\le j\le m .
\end{equation}

\begin{lemma}\label{lem:chain-nodes}
The weights $a_t$ and nodes $w_j$ satisfy \eqref{eq:weight-ranges}.
For every root of unity $\xi\in E$ and every $j$,
$\prod_{\rho\in\Phi}\rho(\xi)^{w_j(\rho)}=1$.
\end{lemma}

\begin{proof}
We have $0\le a_t/b-x\le 1/(kx^{12})$. Both $b$ and $b'$ grow faster
than every fixed power of $x$, while $mh_0$ grows at most polynomially.
At a thick place $b\le w_j\le b+2+mh_0\le b(1+x^{-2})$. At every
other place, including the top places,
$w_j\le(1+x^{-1}+x^{-2})b'+1+mh_0\le2b'$.

Here the only roots of unity in $E$ are $\pm1$. Indeed, an odd prime
$\ell$ dividing the order of a root of unity would ramify in $E$,
so $\ell=q$ or $d$. Both primes exceed $3$ at large scales. Each has
inertia of order two, whereas $\mathbb Q(\zeta_\ell)$ has ramification
index $\ell-1$ at $\ell$. Also, $\mathbb Q(i)$ is unramified at both
$q$ and $d$. If $i$ belonged to $E$, the two inertia groups would fix
it. Their common fixed field is the totally real field $K_e$ of
\S\ref{sec:auxiliary-field}. Finally $|\Phi|=2n$ is even, so
\eqref{eq:chain-nodes} gives an even sum $\sum_\rho w_j(\rho)$ for
every $j$.
\end{proof}

The definite group at level $j$ has rank $N_j=k+j(t_r+1)$ for
$0\le j\le m$. At each $\rho\in\Phi$ take a dominant integral weight
$\lambda^{(0)}_\rho$ of length $k$. At $(0,0,t)$ its first and last
entries are $a_t$ and $-b_t$. Bound the other entries there, and all
entries away from the top places, in absolute value by some
$C_0\le\min_\rho w_1(\rho)$. The $j$th lift uses the weight
\begin{equation}\label{eq:chain-weights}
  \lambda^{(j)}_\rho=
  \bigl(\lambda^{(j-1)}_\rho,
       \underbrace{-w_j(\rho),\ldots,-w_j(\rho)}_{t_r+1}\bigr)
       \qquad(1\le j\le m).
\end{equation}
At a top place $w_1(0,0,t)\ge b_t$, and \eqref{eq:chain-nodes}
increases with $j$. Thus each weight in \eqref{eq:chain-weights} is
dominant. These are the definite ranks and weights of
\citep[\S11.1.2]{adriazola-grh}. For $g=(1,i,h)$, the node in
\eqref{eq:weight-excess} is $w_j(p_tg)=w_j(0,i+1,h+t)$.

\subsection{The lift and the period map}

Give $E^{N_j}$ its standard positive definite hermitian form and write
$U_j=\{\gamma\in\mathrm{GL}_{N_j}(E):\gamma^*\gamma=1\}$.
For fixed finite level, let $\mathcal L_j$ denote the self-dual
$\mathcal O_E$-lattices with level structure, and write $(V_j,\pi_j)$
for the coefficient representation of weight
$(\lambda^{(j)}_\rho)_{\rho\in\Phi}$. A form of level $j$ is a
function $F\colon\mathcal L_j\to V_j$ with
$F(\gamma L)=\pi_j(\gamma)F(L)$ for $\gamma\in U_j$.
Write $\mathcal A_j$ for the space of these forms.

For $L\in\mathcal L_{j-1}$ and $L'\in\mathcal L_j$, the complex
variables of the theta kernel lie in
$\prod_{\rho\in\Phi}M_{N_j\times N_{j-1}}(\mathbb C)$.
Write $K=\mathbb Q(\sqrt{-d})$. Choose a period lattice
$\delta\mathfrak c$ of a CM elliptic curve, with $\mathfrak c$ an
ideal of $\mathcal O_K$. Factor the relative different as
$\mathfrak d_{E/K}=\mathfrak b\bar{\mathfrak b}$ using
\citep[Lemma~7.1]{adriazola-grh}. The period lattice for the two
hermitian lattices is
\[
 \left\{\delta\bigl(\rho(z)\bigr)_{\rho\in\Phi}:
 z\in\mathfrak c\mathfrak b^{-1}
           \operatorname{Hom}_{\mathcal O_E}(L,L')\right\}.
\]
The level data specify a finite linear combination of Heisenberg
translates of the theta function. In the polynomial model, substitute
$X=YD$ in the input polynomial times $\prod_\rho\det(X_\rho)^{n_\rho}$,
replace each monomial in $D$ by the corresponding normalized derivative of the
kernel at zero, and project to $V_j$. With the frame factors of
\citep[Eq.~(7.4)]{adriazola-grh}, denote the resulting map by
$A_j(L',L)\colon V_{j-1}\to V_j$.

\begin{proposition}\label{prop:chain-lift}
Take compatible level data and a Hecke character of infinity type
$(e_\rho)$ trivial on $E^\times$. Suppose the source and target
determinant twists are $\det^{-d_\rho}$ and $\det^{-d'_\rho}$, where
$d_\rho=n_\rho+N_j/2-e_\rho$ and $d'_\rho=N_{j-1}/2-e_\rho$ at
each $\rho\in\Phi$. The class sum
\begin{equation}\label{eq:chain-class-sum}
 (K_jF)(L')=
 \sum_{[L]\in U_{j-1}\backslash\mathcal L_{j-1}}
 \frac{A_j(L',L)F(L)}{|\operatorname{Stab}_{U_{j-1}}(L)|}
\end{equation}
defines a linear map $K_j\colon\mathcal A_{j-1}\to\mathcal A_j$.
\end{proposition}

\begin{proof}
The definite group at fixed finite level has finitely many lattice
classes and finite stabilizers. Both ranks are even. The two
determinant identities give the covariance formula
\citep[Proposition~7.9(iii)]{adriazola-grh}
\[
 A_j(\gamma'L',\gamma L)
    =\pi_j(\gamma')\,A_j(L',L)\,\pi_{j-1}(\gamma)^{-1}.
\]
Replacing $L$ by $\gamma L$ cancels the last factor against the action
on $F(L)$ and conjugates the stabilizer. Replacing $L'$ by $\gamma'L'$
multiplies the sum by $\pi_j(\gamma')$.
\end{proof}

For the arithmetic application the form spaces and lifts must admit
models over a number field. For each auxiliary word $u$, write
$K_{j,u}$ for the lift with its level data and $P_u$ for the final
pairing with the base forms. On the source lattice $N_u$ the required
factorization is \citep[Eq.~(15.5)]{adriazola-grh}
\[
    B_u=\left.(P_u\circ K_{m,u}\circ\cdots\circ K_{1,u})\right|_{N_u},
    \qquad B_u\colon N_u\to R^\vee .
\]
The open arithmetic step is to realize these nonzero maps on the
integral lattices of \S\ref{sec:accounting} and prove its seven
inequalities with the total error of \S\ref{sec:arithmetic-bound}.
Every word must use the chosen $T,u_0$ and the nodes
\eqref{eq:chain-nodes}.
